\chapter{Tanda-tanda Baligh}

\begin{matanbox}[Matan]
Tanda-tanda baligh ada tiga:
\begin{enumerate}
    \item Mencapai genap umur 15 tahun bagi laki–laki dan perempuan.
    \item \textit{Ihtilam} (mimpi basah) pada umur 9 tahun bagi laki–laki dan perempuan.
    \item Haid pada usia 9 tahun bagi perempuan.
\end{enumerate}
\end{matanbox}

\textbf{Tanda-tanda (\textarabic{اَلْعَلَامَاتُ}):} Merupakan bentuk jamak dari kata \textit{'alamah} (\textarabic{عَلَامَة}) artinya petunjuk.

\textbf{Tanda (\textarabic{عَلَامَة}) menurut istilah syariat adalah:} Suatu yang dari keberadaannya berkonsekuensi adanya hukum, namun dari ketiadaannya tidak berkonsekuensi ketiadaan hukum.

\textbf{Maksudnya:} Adanya \textit{ihtilam} berkonsekuensi adanya baligh, namun bukan berarti jika tidak adanya \textit{ihtilam} mesti berkonsekuensi tidak adanya baligh. Karena bisa jadi dia baligh bukan karena tanda \textit{ihtilam}; bisa jadi baligh karena tanda haid, dan bisa jadi karena tanda genap berusia lima belas tahun.

\textbf{Bulugh atau baligh}: Telah sampainya pada batasan \textit{taklif} (pembebanan syariat), dan disyaratkan berakal bersamanya.

\section{Mencapai genap umur 15 tahun bagi laki–laki dan perempuan}

\textbf{Maksudnya:} Berlalunya 15 tahun perhitungan kalendar \textit{qamariyah tahdidiyah} atas \textit{shabiyi} atau \textit{shabiyah} sejak kelahiran sempurnanya.

\textbf{Makna shabiyi:} Usia anak laki-laki dari lahir sampai baligh.

\textbf{Makna shabiyah:} Usia anak perempuan dari lahir sampai baligh.

\textbf{Makna qamariyah tahdidiyah}: Penetapan usia secara tepat dengan menggunakan perhitungan kalender qamariyah yang ditentukan secara tepat.

\section{Ihtilam pada umur 9 tahun bagi laki–laki dan perempuan}

\textbf{Ihtilam:} Apa yang dimimpikan oleh seseorang dalam tidur.
\textbf{Maksudnya:} Keluarnya air mani setelah sempurna berusia 9 tahun dengan perhitungan kalender \textit{qamariyah taqribiyah}.
\textbf{Makna Qamariyah Taqribiyah:} Penetapan usia secara perkiraan dengan menggunakan perhitungan kalender qamariyah tanpa terikat hitungan hari secara persis.

\section{Haid pada usia 9 tahun bagi perempuan}

\textbf{Haid secara istilah syariat:} Darah alami yang keluar dari bagian pangkal terdalam rahim perempuan, dalam keadaan sehat, pada waktu tertentu.

\textbf{Maksudnya:} Perempuan mengalami haid setelah mencapai usia 9 tahun melalui perhitungan kalender \textit{qamariah tahdidiyah}, tidaklah mengapa jika saat haid usianya belum tepat 9 tahun, sekurangnya adalah waktu yang tidak cukup untuk masa haid dan suci; yaitu tidak kurang dari 16 hari. 